\section{Proof of the minimality theorem}\label{app:minimality}\label{sec:minimality}

This appendix shows that no network with four or fewer physical edges can have
equal one-sided slopes at a nested breakpoint, whether the minimum-route
regime grows (an expansion) or shrinks (a contraction) as demand increases.
Together with the five-edge examples this shows that five edges are the
minimum for both kinds of breakpoint. The proof is a finite case analysis.
Each case is closed by an explicit inequality, and Table~\ref{tab:minimality}
lists them all.

\subsection{Setting and statement}

We first fix the class of networks and the kind of breakpoint we study.

\begin{definition}[Standing class]\label{def:standing}
A \emph{network in the standing class} consists of:
\begin{enumerate}
\item a finite directed multigraph $\Gamma$ with vertex set $V$ and edge set
 $\mathcal E$ (parallel and antiparallel edges are allowed);
\item for each edge $e$, a cost $c_e(z)=a_e+b_ez^p$ for $z\ge0$, with
 $a_e,b_e>0$ and one real exponent $p>1$ common to all edges;
\item finitely many commodities $k$, each with an origin $o_k$, a destination
 $d_k\ne o_k$, and demand $tr_k$, where $r_k>0$ and $t>0$ is the scale
 parameter.
\end{enumerate}
The routes of commodity $k$ are \emph{all} simple directed $o_k$--$d_k$
paths in $\Gamma$; we assume there is at least one. We identify a route $P$
with its edge-incidence vector $1_P\in\{0,1\}^{\mathcal E}$. A route is
determined by its edge set, since it is simple.
\end{definition}

A Wardrop equilibrium at scale $t$ is a nonnegative route flow meeting the
demands in which every route with positive flow has minimum cost among the
routes of its commodity. Because every $c_e$ is strictly increasing, the
aggregate edge load $z(t)\in\R^{\mathcal E}_{\ge0}$ of an equilibrium is
unique. Its equilibrium total cost is
\[
 E(t)=\sum_{e}z_e(t)\,c_e(z_e(t)),
\]
and $S(t)$ denotes the minimum of the same total cost over all feasible
route flows. The price of anarchy is $P(t)=E(t)/S(t)$.

We need the set of cheapest routes, not merely the set of used routes.

\begin{definition}[Regimes and nested breakpoints]\label{def:regimeB}
For a load vector $z$, the \emph{complete minimum-route set} $R_k(z)$ of
commodity $k$ is the set of all routes of $k$ whose cost under $c(z)$ is
minimal. It depends only on $z$, not on any route decomposition of $z$. The
\emph{regime} at scale $t$ is the family $R(t)=(R_k(z(t)))_k$. For families we
write $R\subseteq R'$ if $R_k\subseteq R'_k$ for every $k$, and
$R\subsetneq R'$ if moreover the inclusion is strict for some $k$.

A scale $t_0>0$ is a \emph{nested breakpoint} if there is $\varepsilon>0$ such
that $R(t)=R_-$ for $t\in(t_0-\varepsilon,t_0)$ and $R(t)=R_+$ for
$t\in(t_0,t_0+\varepsilon)$, and either
\begin{itemize}
\item $R_-\subsetneq R_+$ (an \emph{expansion}), or
\item $R_+\subsetneq R_-$ (a \emph{contraction}).
\end{itemize}
In either case $R_s$ denotes the smaller and $R_l$ the larger of the two
regimes. The \emph{smaller side} is the punctured one-sided neighbourhood on
which $R(t)=R_s$: it is the left side for an expansion and the right side for
a contraction. The other side is the \emph{larger side}.
\end{definition}

For a family $R$ of route sets let
\[
 \mathcal K_R=\operatorname{span}\{1_P-1_Q:\ P,Q\in R_k\text{ for some }k\}.
\]
This is the set of aggregate edge images of signed route flows on $R$ that
leave every commodity's total demand unchanged. Every element of
$\mathcal K_R$ has zero net flow at every vertex, so $\mathcal K_R$ lies in
the incidence kernel $\mathcal K(\Gamma)$ of $\Gamma$. For a weakly connected
multigraph with $n$ vertices and $e$ edges, $\dim\mathcal K(\Gamma)=e-n+1$.

\begin{theorem}[Four-edge lower bound]\label{thm:fouredge}
Let a network in the standing class have at most four physical edges, and
let $t_0>0$ be a nested breakpoint. Then both one-sided derivatives of $E$
and of $P$ exist at $t_0$, and
\begin{align*}
 E'(t_0^-)&>E'(t_0^+),\quad P'(t_0^-)>P'(t_0^+)
 &&\text{for an expansion},\\
 E'(t_0^-)&<E'(t_0^+),\quad P'(t_0^-)<P'(t_0^+)
 &&\text{for a contraction}.
\end{align*}
\end{theorem}

\begin{corollary}[Minimality of five edges]\label{cor:fiveminimal}
In the standing class, five is the minimum number of physical edges for
which equal one-sided slopes of $E$ (equivalently of $P$) occur at a nested
expansion, and five is also the minimum for a nested contraction.
\end{corollary}

\begin{proof}
Theorem~\ref{thm:fouredge} excludes four or fewer edges. Five edges suffice
for expansions by Theorem~\ref{thm:main}, and for contractions by the
bank-swapped five-edge contraction of Section~\ref{sec:fiveedge}.
\end{proof}

\subsection{General tools}

We use the weak slope comparison, Theorem~\ref{thm:weakcomparison}, and the tangency lemma, Lemma~\ref{lem:tangencycomparison}. Costs in the standing class satisfy their hypotheses: $c_e'(z)=pb_ez^{p-1}$ is continuous on $[0,\infty)$ and positive for $z>0$, and $c_e\ge a_e>0$. For an expansion the smaller side is the left one, so \eqref{eq:distancejump} reads $E'(t_0^-)\ge E'(t_0^+)$; for a contraction it reads $E'(t_0^+)\ge E'(t_0^-)$. Theorem~\ref{thm:fouredge} for $E$ is thus the assertion that the equality $Dv_s=Dv_l$ is impossible. We call this equality \emph{slope equality}. By Lemma~\ref{lem:tangencycomparison}, if $P\in R_{l,k}\setminus R_{s,k}$ and $Q\in R_{s,k}$, then $g(t)=c(z(t))^T(1_P-1_Q)$ is positive on the smaller side, vanishes at $t_0$, and under slope equality has zero one-sided derivative there from the smaller side.

Two symmetries let us halve several case distinctions. The first reverses
all traffic.

\begin{lemma}[Reversal symmetry]\label{lem:reversal}
Let $\Gamma^{\mathrm{rev}}$ be obtained by reversing every edge, with the same
costs, and give each commodity the origin $d_k$ and destination $o_k$ and the
same demand. Then $P\mapsto P^{\mathrm{rev}}$ (traverse in the opposite
direction) is a cost-preserving bijection between the routes of $k$ in the
two networks, with $1_{P^{\mathrm{rev}}}=1_P$. Consequently both networks
have the same equilibrium edge loads at every scale, corresponding regimes,
the same spaces $\mathcal K_R$, and the same functions $E$ and $S$. Nested
breakpoints, their type (expansion or contraction) and their one-sided
slopes correspond.
\end{lemma}

\begin{proof}
A sequence of edges is a simple directed $o_k$--$d_k$ path in $\Gamma$
exactly when the reversed sequence is a simple directed $d_k$--$o_k$ path in
$\Gamma^{\mathrm{rev}}$, and both use the same edges. Hence the incidence
vectors of the route sets agree, so the sets of feasible aggregate loads
agree and route costs agree under every load vector. The Wardrop conditions
and the complete minimum-route sets are statements about these costs, and
$E$ and $S$ depend only on loads. All claims follow.
\end{proof}

The second symmetry is relabelling: renaming vertices and edges by an
isomorphism of multigraphs, with costs and commodities carried along, changes
nothing. We use the composite of reversal with a relabelling, which we call
the \emph{exchange symmetry}; it is described in
Lemma~\ref{lem:sides}\,(e).

Next we reduce to connected graphs and to one commodity per endpoint pair.

\begin{lemma}[Reductions]\label{lem:reductions}
In proving Theorem~\ref{thm:fouredge} we may assume that $\Gamma$ has no
loops, is weakly connected, and has at most one commodity for each ordered
pair of endpoints.
\end{lemma}

\begin{proof}
A loop lies on no simple path; it carries no load and can be deleted, which
does not increase the edge count. If $\Gamma$ has several weak components,
each commodity lives in one of them, and the equilibrium, the regime, $E$ and
$\mathcal K_R$ split as direct sums over components. Every component inherits
constant regimes on both punctured sides, nested in the same direction. On a
component whose regime does not change, the weak comparison theorem with
$R_s=R_l$ gives equal one-sided slopes. Every component whose regime changes
has at most four edges and a nested breakpoint of the same type. If the
inequality of Theorem~\ref{thm:fouredge} is strict for each such component,
adding the components gives strictness for $E$. Finally, commodities with
the same endpoints have the same routes and the same minimum cost; merging
them into one commodity with the summed demand changes neither the
equilibrium loads nor the regimes, nor $E$ or $S$.
\end{proof}

The dimension lemma disposes of every graph whose incidence kernel is at most
one-dimensional.

\begin{lemma}[Dimension lemma]\label{lem:dimension}
At a nested breakpoint of either type, $\mathcal K_{R_s}\subsetneq
\mathcal K_{R_l}$. If $\mathcal K_{R_s}=\{0\}$, slope equality fails, on any
graph in the standing class.
\end{lemma}

\begin{proof}
Choose $k$, $P\in R_{l,k}\setminus R_{s,k}$ and $Q\in R_{s,k}$, and put
$\kappa=1_P-1_Q\in\mathcal K_{R_l}$. On the smaller side, all routes in
$R_{s,k}$ have equal cost, so $c(z(t))$ is orthogonal to $\mathcal K_{R_s}$.
If $\kappa\in\mathcal K_{R_s}$, then $P$ would cost the same as $Q$ there and
would be cheapest, which is false. Hence $\kappa\notin\mathcal K_{R_s}$.

Now let $\mathcal K_{R_s}=\{0\}$. Then all routes in $R_{s,k}$ have the same
incidence vector, so each commodity has a single cheapest route $Q_k$ on the
smaller side and carries its whole demand there. Thus $z(t)=tq$ on the
smaller side, with $q=\sum_kr_k1_{Q_k}=z_0/t_0$, and $q+\mathcal K_{R_s}=\{q\}$
gives $v_s=q$. Suppose slope equality holds. Lemma~\ref{lem:tangencycomparison} gives
$(Dv_s)^T\kappa=(Dq)^T\kappa=0$, and $c_0^T\kappa=0$ because $g(t_0)=0$. For
every edge, $c_e'(z)z=pb_ez^p=p(c_e(z)-a_e)$, also at $z=0$. Hence
\[
 Dq=\frac p{t_0}(c_0-a),
\]
so $(c_0-a)^T\kappa=0$ and then $a^T\kappa=0$. Since
$c_e(tq_e)=a_e+(t/t_0)^p(c_{0,e}-a_e)$,
\[
 c(tq)^T\kappa=a^T\kappa+(t/t_0)^p(c_0-a)^T\kappa=0
\]
for all $t$. On the smaller side $P$ would therefore be cheapest, a
contradiction.
\end{proof}

The parallel-bank lemmas treat a set of parallel edges carrying a fixed total
flow.

\begin{definition}[Banks]\label{def:bank}
A \emph{bank} is a set of $m\ge1$ parallel edges with a common tail and head
and costs $c_i=a_i+b_iz^p$. For a total flow $T\ge0$, the bank equilibrium
assigns loads $x_i\ge0$ with $\sum_ix_i=T$ such that all loaded edges have a
common cost $\ell(T)$, the \emph{effective latency}, and every unloaded edge
has free-flow cost at least $\ell(T)$. The \emph{minimum set} is
$M(T)=\{i:\ a_i\le\ell(T)\}$, the set of edges whose cost equals the minimum.
\end{definition}

The bank equilibrium is explicit: with
$N(\lambda)=\sum_i\bigl((\lambda-a_i)_+/b_i\bigr)^{1/p}$, which is continuous,
zero at $\lambda=\min_ia_i$ and strictly increasing for larger $\lambda$,
$\ell(T)$ is the unique $\lambda\ge\min_ia_i$ with $N(\lambda)=T$, and
$x_i=((\ell-a_i)_+/b_i)^{1/p}$. The next lemma records how the latency
behaves near a given flow.

\begin{lemma}[Single bank]\label{lem:bank}
Let $T_0>0$, $\lambda_0=\ell(T_0)$, $I=\{i:a_i<\lambda_0\}$ (the positively
loaded edges, nonempty) and $J=\{j:a_j=\lambda_0\}$. Put
$d_i=c_i'(x_i)>0$ for $i\in I$ and $\eta=\sum_{i\in I}d_i^{-1}$.
\begin{enumerate}
\item[(a)] $\ell$ is continuous and strictly increasing, so $M(T)$ is
 nondecreasing in $T$ for inclusion: branches only join as flow grows.
\item[(b)] $\ell$ has left derivative $\ell'(T_0^-)=\eta^{-1}>0$.
\item[(c)] If $J=\emptyset$, then $\ell$ is real-analytic near $T_0$ with
 $\ell'(T_0)=\eta^{-1}$. If $J\ne\emptyset$, the right derivative is
 $\ell'(T_0^+)=0$.
\end{enumerate}
\end{lemma}

\begin{proof}
(a) $N$ is a continuous strictly increasing bijection from
$[\min_ia_i,\infty)$ onto $[0,\infty)$, so its inverse $\ell$ is continuous
and strictly increasing, and $\{i:a_i\le\ell(T)\}$ grows with $T$.

(b), (c) Near $\lambda_0$, $N(\lambda)=N_I(\lambda)+N_J(\lambda)$ with
$N_I(\lambda)=\sum_{i\in I}((\lambda-a_i)/b_i)^{1/p}$ real-analytic near
$\lambda_0$, $N_I'(\lambda_0)=\eta$ (because $dx_i/d\lambda=1/c_i'(x_i)$), and
$N_J(\lambda)=\sum_{j\in J}((\lambda-\lambda_0)_+/b_j)^{1/p}$, which vanishes
for $\lambda\le\lambda_0$. Hence $\ell=N_I^{-1}$ to the left of $T_0$, which
gives (b), and also to the right when $J=\emptyset$, which gives the first
part of (c). If $J\ne\emptyset$, put $C=\sum_{j\in J}b_j^{-1/p}>0$. For
$T>T_0$ and $\lambda=\ell(T)>\lambda_0$, monotonicity of $N_I$ gives
$T-T_0\ge C(\lambda-\lambda_0)^{1/p}$, so
$0<\ell(T)-\lambda_0\le((T-T_0)/C)^p=o(T-T_0)$, since $p>1$.
\end{proof}

The curvature lemma compares the effective latency of a bank with a single
common-power edge.

\begin{lemma}[Parallel-bank curvature]\label{lem:curvature}
Let $\ell$ be the effective latency of a bank and let $T>0$ be such that all
$m$ edges are positively loaded, $x_i>0$. Put $d_i=c_i'(x_i)$ and
$\eta=\sum_id_i^{-1}$. Then
\begin{equation}\label{eq:min_curv}
 \ell''(T)-\frac{p-1}{T}\ell'(T)
 =\frac{p-1}{T\eta^3}\Bigl[T\sum_i\frac1{x_id_i^2}-\eta^2\Bigr]\ge0,
\end{equation}
with equality exactly when all offsets $a_i$ are equal. In that case
$\ell(T)=a+\kappa T^p$ for all $T\ge0$, with $a$ the common offset and
$\kappa=(\sum_ib_i^{-1/p})^{-p}$.
\end{lemma}

\begin{proof}
With all edges loaded, $T=N(\lambda)$ with
$N'(\lambda)=\sum_i1/c_i'(x_i)=\eta$ and
$N''(\lambda)=-\sum_ic_i''(x_i)/c_i'(x_i)^3$. For common-power costs
$c_i''(x)=(p-1)c_i'(x)/x$ when $x>0$, so
$N''=-(p-1)\sum_i1/(x_id_i^2)$. Inverse differentiation gives
$\ell'=1/N'=\eta^{-1}$ and
\[
 \ell''=-\frac{N''}{(N')^3}=\frac{p-1}{\eta^3}\sum_i\frac1{x_id_i^2}.
\]
Subtracting $(p-1)\ell'/T=(p-1)\eta^2/(T\eta^3)$ gives the identity in
\eqref{eq:min_curv}. By the Cauchy--Schwarz inequality,
\[
 \eta^2=\Bigl(\sum_i\sqrt{x_i}\cdot\frac1{\sqrt{x_i}\,d_i}\Bigr)^2
 \le\Bigl(\sum_ix_i\Bigr)\sum_i\frac1{x_id_i^2}=T\sum_i\frac1{x_id_i^2},
\]
with equality exactly when $x_i d_i$ is the same for all $i$. Since
$x_id_i=pb_ix_i^p=p(\ell(T)-a_i)$, this means all $a_i$ are equal. If
$a_i=a$ for all $i$, then $x_i=((\ell-a)/b_i)^{1/p}$ and summing gives
$T=(\ell-a)^{1/p}\sum_ib_i^{-1/p}$, which is the stated formula.
\end{proof}

We use the lemma through the operator $\mathcal L[\varphi](t)=\varphi''(t)-
(p-1)\varphi'(t)/t$ on functions of $t>0$. It annihilates $\alpha+\beta t^p$
for all constants $\alpha,\beta$. If $Q>0$ and all bank edges are positively
loaded at $T=Qt$, the function $\ell_Q(t)=\ell(Qt)$ satisfies
\begin{equation}\label{eq:min_tcurv}
 \mathcal L[\ell_Q](t)=Q^2\Bigl[\ell''(Qt)-\frac{p-1}{Qt}\ell'(Qt)\Bigr]>0
 \quad\text{unless all offsets are equal.}
\end{equation}
If $\varphi(t_0)=\varphi'(t_0)=0$, then $\varphi''(t_0)=\mathcal L[\varphi](t_0)$.

The next lemma disposes of every network whose underlying graph is a tree.

\begin{lemma}[Trees of parallel banks]\label{lem:tree}
Suppose that the underlying simple undirected graph of $\Gamma$ (obtained by
merging parallel and antiparallel edges) is a tree. Then no nested
contraction occurs, and at every nested expansion
$E'(t_0^-)>E'(t_0^+)$.
\end{lemma}

\begin{proof}
For an ordered pair $(x,y)$ of adjacent vertices, the edges from $x$ to $y$
form a bank $\mathcal B_{xy}$. Every commodity $k$ has a unique vertex
sequence, namely the tree path from $o_k$ to $d_k$, and its routes are
exactly the choices of one edge from each bank along that sequence; all such
choices are simple paths and hence available. The bank $\mathcal B_{xy}$
therefore carries the total flow $Q_{xy}t$, where $Q_{xy}$ is the sum of the
$r_k$ over commodities whose sequence traverses $x$ then $y$. The
equilibrium splits each bank independently as a bank equilibrium, because a
route's cost is the sum of the costs of its bank edges. A route is cheapest
exactly when each of its edges lies in the minimum set of its bank. Thus
$R_k(t)$ is the product of the sets $M_{xy}(Q_{xy}t)$ along the sequence of
$k$, and the regime is nondecreasing in $t$ by Lemma~\ref{lem:bank}(a). This
excludes $R_+\subsetneq R_-$.

At an expansion, every bank with $Q_{xy}>0$ has minimum sets $M_-\subseteq
M_+$ on the two sides (the product of nonempty sets is monotone in each
factor), and at least one such bank has $M_-\subsetneq M_+$. Fix a bank with
$Q=Q_{xy}>0$ and $M_-\subsetneq M_+$, and write $T_0=Qt_0$ and
$\lambda_0=\ell(T_0)$. If $i\in M_-$ then $a_i\le\ell(Qt)<\lambda_0$ for
$t<t_0$; if $j\in M_+\setminus M_-$ then $a_j>\ell(Qt)$ for $t<t_0$ and
$a_j\le\ell(Qt)$ for $t>t_0$, so $a_j=\lambda_0$. In the notation of
Lemma~\ref{lem:bank}, $I=M_-$ and $J=M_+\setminus M_-\ne\emptyset$. The bank
contributes $E_{xy}(t)=Qt\,\ell(Qt)$ to $E$, so by Lemma~\ref{lem:bank}(b,c)
\[
 E_{xy}'(t_0^-)-E_{xy}'(t_0^+)=Q^2t_0\bigl[\ell'(T_0^-)-\ell'(T_0^+)\bigr]
 =Q^2t_0\eta^{-1}>0.
\]
For a bank whose minimum set $M$ does not change, the same argument shows
$a_i<\lambda_0$ for $i\in M$ and $a_i>\lambda_0$ otherwise, so $J=\emptyset$
and its contribution has a continuous derivative by
Lemma~\ref{lem:bank}(c). Summing over banks gives the claim.
\end{proof}

We also need that a single commodity's cost increases with its demand.

\begin{lemma}[Monotone single-commodity cost]\label{lem:monotone}
If all commodities share the same endpoints and are merged into one
commodity of demand $Rt$, its equilibrium cost $\lambda(t)$ is strictly
increasing in $t$.
\end{lemma}

\begin{proof}
Let $t_1<t_2$, with loads $z_i=z(t_i)$ and costs $\lambda_i$. Every route
costs at least $\lambda_i$ under $c(z_i)$, and the used routes cost exactly
$\lambda_i$. Decomposing $z_j$ into routes therefore gives
$c(z_i)^Tz_j\ge\lambda_iRt_j$ and $c(z_i)^Tz_i=\lambda_iRt_i$. Hence
\[
 0\le(c(z_2)-c(z_1))^T(z_2-z_1)\le(\lambda_2-\lambda_1)R(t_2-t_1).
\]
If $\lambda_2\le\lambda_1$, the middle expression vanishes, so $z_1=z_2$ by
strict monotonicity of the costs. This is impossible, because the total flow
leaving the origin is $Rt_1$ for one load and $Rt_2$ for the other.
\end{proof}

\subsection{The transitive triangle with a duplicated side}

After the general tools, only one family of graphs remains. The next lemma
identifies it.

\begin{lemma}[Graph reduction]\label{lem:graphs}
Let $\Gamma$ be loopless and weakly connected with $n$ vertices and $e\le4$
edges, and suppose its underlying simple graph is not a tree. If
$\dim\mathcal K(\Gamma)\le1$, a nested breakpoint has
$\mathcal K_{R_s}=\{0\}$. Otherwise $n=3$, $e=4$, and if
$\dim\mathcal K_{R_s}\ge1$ then $\mathcal K_{R_l}=\mathcal K(\Gamma)$ and
$\Gamma$ is a \emph{transitive triangle with one duplicated side}: up to
relabelling, its vertices are $s,u,v$, it has one edge on each of the sides
$s\to v$, $s\to u$, $u\to v$, and one further edge parallel to one of them.
\end{lemma}

\begin{proof}
$\dim\mathcal K(\Gamma)=e-n+1$. If $n\ge4$, this is at most $1$. If $n=2$,
the underlying simple graph is a single edge, a tree. If $n=3$ and $e\le3$, it
is at most $1$. When it is at most one, the strict inclusion
$\mathcal K_{R_s}\subsetneq\mathcal K_{R_l}\subseteq\mathcal K(\Gamma)$ of
Lemma~\ref{lem:dimension} forces $\mathcal K_{R_s}=\{0\}$. Hence the only
remaining case is $n=3$, $e=4$, $\dim\mathcal K(\Gamma)=2$. If
$\dim\mathcal K_{R_s}\ge1$, the strict inclusion forces
$\mathcal K_{R_l}=\mathcal K(\Gamma)$.

The cost vector $c_0$ is orthogonal to $\mathcal K_{R_l}$, being the limit of
cost vectors orthogonal to it on the larger side. The incidence vector of a
directed cycle lies in $\mathcal K(\Gamma)=\mathcal K_{R_l}$, but its inner
product with $c_0>0$ is positive. Hence $\Gamma$ has no directed cycle, and
in particular no pair of antiparallel edges. An acyclic graph on three
vertices has a topological order $s<u<v$ with every edge going forward. The
underlying simple graph is not a tree, so all three pairs are joined; their
multiplicities are positive and sum to four, so they are $2,1,1$ in some
order.
\end{proof}

For the rest of this subsection $\Gamma$ is a transitive triangle with one
duplicated side. Its \emph{sides} are $f$ ($s\to v$), $h$ ($s\to u$) and $g$
($u\to v$). The duplicated side $\sigma^*$ has two edges $e_1,e_2$; each other
side has one edge, also denoted $f$, $h$ or $g$. Only the endpoint pairs
$(s,v)$, $(s,u)$, $(u,v)$ are joined by directed paths. By
Lemma~\ref{lem:reductions} there is at most one commodity per pair; we write
$R,A,B\ge0$ for the demand coefficients of $s\to v$, $s\to u$, $u\to v$,
with a zero coefficient meaning that the commodity is absent. The routes are
single edges of side $h$ (for $s\to u$) or $g$ (for $u\to v$); and for
$s\to v$ the \emph{direct} routes (one edge of side $f$) and the
\emph{indirect} routes (one edge of side $h$ followed by one edge of side
$g$).

The next lemma describes regimes through side minima; it is where
availability of all simple paths is used.

\begin{lemma}[Side structure and aggregation]\label{lem:sides}
Fix a load vector $z$. For each side $\sigma$ let $m_\sigma$ be the least
cost among its edges and $M_\sigma$ the set of edges attaining it. Put
$\varpi=m_h+m_g-m_f$.
\begin{enumerate}
\item[(a)] $R_{su}=M_h$ and $R_{uv}=M_g$ (when present), and
 $R_{sv}=M_f$ if $\varpi>0$, $R_{sv}=M_h\times M_g$ if $\varpi<0$, and
 $R_{sv}=M_f\cup(M_h\times M_g)$ if $\varpi=0$.
\item[(b)] An edge outside $M_\sigma$ carries zero load in any equilibrium.
\item[(c)] If the sign of $\varpi$ is nonzero and fixed, the side loads
 $y_\sigma=\sum_{e\in\sigma}z_e$ of an equilibrium at scale $t$ are
 $(y_f,y_h,y_g)=(R,A,B)t$ when $\varpi>0$ and $(0,A+R,B+R)t$ when $\varpi<0$.
\item[(d)] If a route of commodity $k$ is not cheapest, its excess cost over
 the cheapest route of $k$ is $m_\sigma'-m_\sigma$ when it uses a
 non-minimum edge of cost $m_\sigma'$ on a side $\sigma$ and otherwise
 follows cheapest edges, and it is $|\varpi|$ when it is a direct (resp.\
 indirect) route built from cheapest edges while $\varpi<0$ (resp.\
 $\varpi>0$). In particular, all such routes of the same kind have the same
 excess, whichever edge of a bank with equal costs they use.
\item[(e)] (Exchange symmetry.) Reversal (Lemma~\ref{lem:reversal}) followed
 by the relabelling $s\leftrightarrow v$ maps the network to a transitive
 triangle with one duplicated side in which the roles of $(h,A)$ and $(g,B)$
 are exchanged and $f$ and $R$ are unchanged; it preserves the type of
 every breakpoint, the regimes' sign of $\varpi$, and the slopes of $E$.
\end{enumerate}
\end{lemma}

\begin{proof}
(a) Every combination of an $h$-edge and a $g$-edge is a simple
$s\to v$ path and is available. The cheapest $s\to u$ and $u\to v$ routes are
therefore the cheapest edges of the corresponding side, the cheapest
indirect routes are the pairs in $M_h\times M_g$ with cost $m_h+m_g$, and the
cheapest direct routes are $M_f$ with cost $m_f$. Comparing $m_f$ with
$m_h+m_g$ gives (a).

(b) If an edge $e$ of side $\sigma$ carries load, some used route contains
it, and replacing $e$ by an edge of $M_\sigma$ gives an available route that
is strictly cheaper unless $e\in M_\sigma$.

(c) If $\varpi>0$, all $s\to v$ flow is direct, and if $\varpi<0$ it is all
indirect; the other commodities have routes inside a single side.

(d) This is the computation in (a).

(e) Reversal turns $s\to u$, $u\to v$, $s\to v$ into $u\to s$, $v\to u$,
$v\to s$. After renaming $s\leftrightarrow v$ these read $u\to v$, $s\to u$,
$s\to v$: the old side $g$ becomes the new side $h$ and conversely, and the
old $s\to u$ commodity becomes the new $u\to v$ commodity. Direct routes stay
direct. Lemma~\ref{lem:reversal} gives the remaining assertions.
\end{proof}

Write $\pi=1_{e_2}-1_{e_1}$ and, for edges $\varphi,\chi,\psi$ of the sides
$f,h,g$, let $\tau(\varphi,\chi,\psi)=1_\chi+1_\psi-1_\varphi$ be a
\emph{triangle difference}. Two triangle differences are equal or differ by
$\pm\pi$. We say that the duplicated side is \emph{active} in a regime if
some route of the regime uses an edge of $\sigma^*$.

The classification lemma lists the possible route-difference spaces.

\begin{lemma}[Route-space classification]\label{lem:routespaces}
$\mathcal K(\Gamma)$ is two-dimensional and is spanned by $\pi$ and any
triangle difference. For any family $R$ of nonempty route sets,
$\mathcal K_R$ is spanned by vectors each of which is $\pm\pi$ or a triangle
difference; hence if $\dim\mathcal K_R=1$, either $\mathcal K_R=
\operatorname{span}\{\pi\}$ or $\mathcal K_R=\operatorname{span}\{\tau\}$ for
a single triangle difference $\tau$. For a complete regime with data
$M_\sigma,\varpi$ as in Lemma~\ref{lem:sides}:
\begin{enumerate}
\item[(a)] $\pi\in\mathcal K_R$ if and only if $\sigma^*$ is active and
 $M_{\sigma^*}=\{e_1,e_2\}$;
\item[(b)] $\mathcal K_R$ contains a triangle difference if and only if
 $R>0$ and $\varpi=0$;
\item[(c)] $\dim\mathcal K_R$ is the number of conditions among (a), (b)
 that hold, and when only (b) holds, $\mathcal K_R=\operatorname{span}\{\tau\}$
 where $\tau$ is built from the unique cheapest edge of each side.
\end{enumerate}
\end{lemma}

\begin{proof}
$\dim\mathcal K(\Gamma)=4-3+1=2$, and $\pi$ and a triangle difference are
independent elements of it, since a triangle difference has three nonzero
entries on three different sides while $\pi$ is supported on one side.
The generators $1_P-1_Q$ of $\mathcal K_R$ are differences of two routes of
one commodity. For $s\to u$ or $u\to v$ they are $0$ or $\pm\pi$. For
$s\to v$, two direct routes or two indirect routes differ by $0$ or $\pm\pi$,
and a direct and an indirect route differ by a triangle difference. If
$\dim\mathcal K_R=1$, either every nonzero generator is $\pm\pi$, or some
generator is a triangle difference $\tau$; in the latter case no generator is
$\pm\pi$ (which is not a multiple of $\tau$), and a second, different
triangle difference would differ from $\tau$ by $\pm\pi$ and raise the
dimension.

For a complete regime, routes use only edges of the sets $M_\sigma$, and by
Lemma~\ref{lem:sides}(a) every combination of such edges compatible with
the sign of $\varpi$ is a route of the regime. If $\sigma^*$ is active and
$|M_{\sigma^*}|=2$, a route through $e_1$ and the same route through $e_2$
both lie in the regime, so $\pi\in\mathcal K_R$. Otherwise all routes of the
regime use a fixed edge on each side, and $\mathcal K_R$ is contained in the
span of the single triangle difference $\tau$ built from these edges, which
does not contain $\pi$. This proves (a). A triangle difference is a
generator exactly when $s\to v$ has both a cheapest direct and a cheapest
indirect route, that is $R>0$ and $\varpi=0$; if this fails, every generator
is $0$ or $\pm\pi$. This proves (b), and (c) follows.
\end{proof}

In the remaining two lemmas $\Gamma$ is a transitive triangle with one
duplicated side, $t_0$ is a nested breakpoint of either type,
$\dim\mathcal K_{R_s}=1$ and $\mathcal K_{R_l}=\mathcal K(\Gamma)$. By
Lemma~\ref{lem:routespaces}, the larger regime has $R>0$, $\varpi=0$, both
$e_1,e_2$ cheapest and $\sigma^*$ active. By the same lemma, the smaller
regime has one of two types:
\begin{description}
\item[Type P] (old cycle is the parallel pair) $\mathcal K_{R_s}=
 \operatorname{span}\{\pi\}$: on the smaller side $\sigma^*$ is active with
 both edges cheapest, and $\varpi\ne0$.
\item[Type T] (old cycle is the triangle) $\mathcal K_{R_s}=
 \operatorname{span}\{\tau\}$: on the smaller side $\varpi=0$, and exactly one
 edge of $\sigma^*$, say $e_1$, is cheapest.
\end{description}
In type~P the sign of $\varpi$ on the smaller side is constant, because
$R_{sv}$ is either a set of direct routes or a set of indirect routes; we
call the \emph{old long route} direct if $\varpi>0$ and indirect if
$\varpi<0$ there.

\begin{lemma}[Old cycle is the parallel pair]\label{lem:typeP}
In type~P, slope equality fails, for expansions and for contractions.
\end{lemma}

\begin{proof}
\emph{Set-up.} By Lemma~\ref{lem:sides}(c), on the smaller side the side
loads are $y_\sigma(t)=Q_\sigma t$, with $(Q_f,Q_h,Q_g)=(R,A,B)$ for a direct
old long route and $(0,A+R,B+R)$ for an indirect one. The bank
$\sigma^*=\{e_1,e_2\}$ is active, so $Q:=Q_{\sigma^*}>0$; let $\ell$ be its
effective latency, $T_0=Qt_0$ and $\lambda_0=\ell(T_0)$. If the old long
route is indirect, side $f$ carries no load and is not on any route of
$R_s$, so $\sigma^*\ne f$. Hence $\sigma^*=f$ forces a direct old long
route. By the exchange symmetry (Lemma~\ref{lem:sides}(e)) we may assume
$\sigma^*\in\{f,g\}$; for $\sigma^*=g$ with a direct old long route, activity
means $B>0$.

For $t$ near $t_0$ define the \emph{gap}
\[
 G(t)=
 \begin{cases}
 c_h(At)+c_g(Bt)-\ell(Rt), & \sigma^*=f\ \text{(direct)},\\
 a_f-c_h((A+R)t)-\ell((B+R)t), & \sigma^*=g,\ \text{indirect},\\
 c_h(At)+\ell(Bt)-c_f(Rt), & \sigma^*=g,\ \text{direct}.
 \end{cases}
\]
On the smaller side, $G(t)$ is the cost of a new route (indirect, direct and
indirect, respectively) minus the cost of an old one; both edges of the bank
have cost $\ell$, so all new routes have this same gap
(Lemma~\ref{lem:sides}(d)). Thus $G>0$ on the smaller side, $G(t_0)=0$, and
under slope equality Lemma~\ref{lem:tangencycomparison} gives $G'(t_0)=0$ as a
one-sided derivative from the smaller side.

\emph{Both bank edges are loaded at $t_0$ for an expansion.} The bank load
on the smaller side is $Qt$, so the bank's common cost there is $\ell(Qt)$,
strictly increasing in $t$. For an expansion the smaller side is the left
one; both edges are cheapest there, so $a_i\le\ell(Qt)<\lambda_0$ for
$t<t_0$, whence $a_1,a_2<\lambda_0$ and both are loaded at $t_0$. For a
contraction we only get $a_i\le\lambda_0$. Call the bank
\emph{nondegenerate} if $a_1,a_2<\lambda_0$ and \emph{degenerate} otherwise;
degeneracy can occur only for contractions, and then exactly one edge, say
$e_2$, has $a_{e_2}=\lambda_0$ and zero load at $t_0$ (the bank load
$T_0>0$ is carried by $e_1$).

\emph{Nondegenerate bank.} By Lemma~\ref{lem:bank}(c), $\ell$ is analytic
near $T_0$ with $\ell'>0$, so $G$ is $C^2$ near $t_0$ and its derivatives
are two-sided.

\smallskip\noindent
\emph{(P1) $\sigma^*=f$.} $G(t)=\alpha+\beta t^p-\ell(Rt)$ with
$\alpha=a_h+a_g$, $\beta=b_hA^p+b_gB^p$. If the offsets of the bank differ,
\eqref{eq:min_tcurv} gives $G''(t_0)=\mathcal L[G](t_0)=-\mathcal
L[\ell_R](t_0)<0$ under $G(t_0)=G'(t_0)=0$. Then $G<0$ on both punctured sides
near $t_0$, contradicting $G>0$ on the smaller side.

\smallskip\noindent
\emph{(P2) $\sigma^*=g$, indirect.}
$G'(t_0)=-pb_h(A+R)^pt_0^{p-1}-(B+R)\ell'((B+R)t_0)<0$, contradicting
$G'(t_0)=0$.

\smallskip\noindent
\emph{(P3) $\sigma^*=g$, direct.} $G(t)=\ell(Bt)-\alpha-\beta t^p$ with
$\alpha=a_f-a_h$ and $\beta=b_fR^p-b_hA^p$. If the offsets differ,
$G''(t_0)=\mathcal L[\ell_B](t_0)>0$ under $G(t_0)=G'(t_0)=0$, so $G>0$ on both
punctured sides of $t_0$ near $t_0$. On the larger side consider the
\emph{old continuation}: $s\to v$ entirely on $f$, $s\to u$ on $h$, and
$u\to v$ split over the bank as its bank equilibrium with total $Bt$, both
edges loaded for $t$ near $t_0$. Every used route is cheapest: the $s\to u$
route is the only one; the two $u\to v$ routes cost $\ell(Bt)$; and every
indirect $s\to v$ route costs $c_h(At)+\ell(Bt)$, which exceeds the direct cost
by $G(t)>0$. So the old continuation is an equilibrium; by uniqueness it is
the equilibrium. Its complete minimum-route set is $R_s$, since the
substituted indirect routes are not cheapest (their gap is $G(t)>0$), even
though $R_l$ contains them. This contradicts $R(t)=R_l$ on the larger side.

\smallskip\noindent
\emph{Equal offsets in (P1) or (P3).} By Lemma~\ref{lem:curvature},
$\ell(T)=a^*+\kappa T^p$, so $G(t)=\alpha'+\beta't^p$ for constants
$\alpha',\beta'$. Then $G'(t_0)=p\beta't_0^{p-1}=0$ gives $\beta'=0$, and
$G(t_0)=0$ gives $\alpha'=0$. Hence $G\equiv0$, contradicting $G>0$ on the
smaller side.

\smallskip\noindent
\emph{Degenerate bank (contractions only; smaller side is the right side).}
Here $\ell'(T_0^+)=0$ by Lemma~\ref{lem:bank}(c). On the larger (left) side
both bank edges are cheapest, so the bank's common cost is at least
$a_{e_2}=\lambda_0$; as the bank's loads form a bank equilibrium for its total
$y_{\sigma^*}(t)$, this means $\ell(y_{\sigma^*}(t))\ge\lambda_0$, that is
$y_{\sigma^*}(t)\ge T_0$ for $t<t_0$. Let $x(t)\ge0$ denote the indirect
$s\to v$ flow on the left; the left side loads are $y_f=Rt-x$, $y_h=At+x$,
$y_g=Bt+x$, and $\varpi=0$ there.

\emph{(D1) $\sigma^*=f$.} $y_f(t)=Rt-x(t)\le Rt<Rt_0=T_0$ for $t<t_0$, a
contradiction.

\emph{(D2) $\sigma^*=g$, indirect.} The right derivative is
$G'(t_0^+)=-pb_h(A+R)^pt_0^{p-1}-(B+R)\ell'(T_0^+)=-pb_h(A+R)^pt_0^{p-1}<0$,
contradicting $G'(t_0^+)=0$.

\emph{(D3) $\sigma^*=g$, direct.} The right derivative of
$G=\ell(Bt)-\alpha-\beta t^p$ is $-p\beta t_0^{p-1}$, so slope equality gives
$\beta=0$, i.e.\ $b_fR^p=b_hA^p$ (hence $A>0$), and then $G(t_0)=0$ gives
$\alpha=\lambda_0$, i.e.\ $a_f=a_h+\lambda_0$. On the left, $\varpi=0$
and the $g$-bank cost $\lambda_g(t)=\ell(y_g(t))\ge\lambda_0$ give
\[
 a_f+b_f(Rt-x)^p=a_h+b_h(At+x)^p+\lambda_g(t)
 \ge a_h+b_hA^pt^p+\lambda_0=a_f+b_fR^pt^p .
\]
Since $0\le Rt-x\le Rt$, equality holds throughout: $x(t)=0$ and
$\lambda_g(t)=\lambda_0$. But then $y_g(t)=Bt<T_0$, so
$\lambda_g(t)=\ell(Bt)<\lambda_0$, a contradiction.

The exchange symmetry transfers (P2), (P3), (D2), (D3) to $\sigma^*=h$, which
completes type~P.
\end{proof}

For type T we parametrise the smaller side by the indirect flow and use the
implicit-function theorem.

\begin{lemma}[Old cycle is the triangle]\label{lem:typeT}
In type~T, slope equality fails, for expansions and for contractions.
\end{lemma}

\begin{proof}
\emph{Set-up.} On the smaller side $e_2$ is not cheapest, so it carries no
load (Lemma~\ref{lem:sides}(b)); it is unused and, more strongly, not
minimum. On the larger side it is cheapest. Since $\varpi=0$ on both sides,
every side lies on some $s\to v$ route of both regimes. Let $Q$ be such a
route through $e_1$ built from cheapest edges of the smaller side, and $P$ the
same route with $e_1$ replaced by $e_2$. By Lemma~\ref{lem:sides}(a),
$Q\in R_{s,sv}$ and $P\in R_{l,sv}\setminus R_{s,sv}$. Their gap is
\[
 \vartheta(t)=a_{e_2}-c_{e_1}(z_{e_1}(t)),
\]
which by Lemma~\ref{lem:tangencycomparison} is positive on the smaller side and zero at
$t_0$. Write $c_\sigma$ for the cost of the single cheapest edge of side
$\sigma$ on the smaller side ($c_{\sigma^*}=c_{e_1}$), and $y_\sigma$ for its
load, so $y_{\sigma^*}=z_{e_1}$.

\emph{The duplicated edge is loaded at $t_0$.} If $y_{\sigma^*}(t_0)=0$, then
$\vartheta(t_0)=0$ gives $a_{e_2}=a_{e_1}$, and on the smaller side
$c_{e_1}(z_{e_1})\ge a_{e_1}=a_{e_2}$, so $e_2$ would be cheapest there. Hence
$y_{\sigma^*}(t_0)>0$, $c_{e_1}'(y_{\sigma^*}(t_0))>0$, and under slope
equality the tangency $\vartheta'(t_0)=0$ (from the smaller side) is equivalent
to
\begin{equation}\label{eq:min_stationary}
 y_{\sigma^*}'(t_0)=0\qquad\text{(one-sided, from the smaller side)}.
\end{equation}

\emph{The smaller-side branch.} On the smaller side $\varpi=0$, so with
$x(t)\ge0$ the indirect $s\to v$ flow,
\[
 y_f=Rt-x,\quad y_h=At+x,\quad y_g=Bt+x,\qquad
 \Psi(x,t):=c_f(y_f)-c_h(y_h)-c_g(y_g)=0 .
\]
Extend each cost to negative loads by $\tilde c(z)=a+b\operatorname{sgn}(z)
|z|^p$, which is $C^1$ with $\tilde c'(z)=pb|z|^{p-1}\ge0$. At $t_0$ the loads
$y_f,y_h$ sum to $(R+A)t_0>0$, so one of the three derivatives
$d_\sigma=\tilde c_\sigma'(y_\sigma(t_0))$ is positive, and
$\partial_x\Psi=-(d_f+d_h+d_g)=:-J<0$. By the implicit-function theorem there
is a unique $C^1$ solution $\tilde x(t)$ of $\Psi=0$ near $(x(t_0),t_0)$;
by continuity of $z$ the actual $x(t)$ on the smaller side coincides with
it. Implicit differentiation gives, as one-sided derivatives from the
smaller side,
\begin{equation}\label{eq:min_velocities}
 y_f'=\frac{(R+A)d_h+(R+B)d_g}J,\
 y_h'=\frac{(R+A)d_f-(B-A)d_g}J,\
 y_g'=\frac{(R+B)d_f-(A-B)d_h}J .
\end{equation}
Recall $d_\sigma>0$ exactly when $y_\sigma(t_0)>0$. By the exchange symmetry
we may assume $\sigma^*\in\{f,h\}$.

\smallskip\noindent
\emph{(T1) $\sigma^*=f$.} By \eqref{eq:min_velocities},
$y_f'(t_0)=0$ forces $d_h=d_g=0$, so $y_h(t_0)=At_0+x(t_0)=0$ and
$y_g(t_0)=Bt_0+x(t_0)=0$, whence $A=B=0$ and $x(t_0)=0$. Then $s\to v$ is the
only commodity, and its equilibrium cost at $t_0$ is the cost of the
(cheapest) indirect route, $\lambda(t_0)=a_h+a_g$. On the left side of $t_0$,
whichever regime holds there, $\varpi=0$, so the indirect route is cheapest
and $\lambda(t)=c_h(y_h(t))+c_g(y_g(t))\ge a_h+a_g=\lambda(t_0)$ for $t<t_0$.
This contradicts Lemma~\ref{lem:monotone}.

\smallskip\noindent
\emph{(T2) $\sigma^*=h$.} Here $y_h(t_0)>0$, so $d_h>0$, and the
inequality $\vartheta>0$ on the smaller side says
\begin{equation}\label{eq:min_hbelow}
 y_h(t)<y_h(t_0)\qquad\text{on the smaller side.}
\end{equation}
We use $L=R+A>0$ and $M=B-A$, so $y_f=Lt-y_h$ and $y_g=Mt+y_h$, and
\eqref{eq:min_velocities} reads $y_h'=(Ld_f-Md_g)/J$.

\emph{(T2a) $y_f(t_0)>0$ and $y_g(t_0)=0$.} Then $y_h'(t_0)=Ld_f/J>0$,
contradicting \eqref{eq:min_stationary}.

\emph{(T2b) $y_f(t_0)=0$.} Then $x(t_0)=Rt_0$, so $y_g(t_0)=(B+R)t_0>0$ and
$y_h'(t_0)=-Md_g/J$. By \eqref{eq:min_stationary}, $M=0$, so $y_g=y_h$ on the
smaller side, and $\Psi=0$ reads $c_f(y_f)=\Theta(y_h)$ with
$\Theta(y)=c_h(y)+c_g(y)$ strictly increasing. Since $y_f(t)\ge0$,
\[
 \Theta(y_h(t))=c_f(y_f(t))\ge a_f=c_f(y_f(t_0))=\Theta(y_h(t_0)),
\]
so $y_h(t)\ge y_h(t_0)$ on the smaller side, contradicting
\eqref{eq:min_hbelow}. (For an expansion one may instead note that
$y_f'(t_0)>0$ by \eqref{eq:min_velocities} would make $y_f$ negative just to
the left; the argument given covers both types.)

\emph{(T2c) $y_f(t_0)>0$ and $y_g(t_0)>0$.} All three loads are positive at
$t_0$, so the costs are analytic near them and the branch $y_h(t)$ is
analytic. Condition \eqref{eq:min_stationary} gives $U:=Ld_f=Md_g>0$, so
$M>0$. Differentiating $c_f(Lt-y_h)=c_h(y_h)+c_g(Mt+y_h)$ twice at a point where
$y_h'=0$ gives $Jy_h''=L^2c_f''-M^2c_g''$. With
$c''(y)=(p-1)c'(y)/y$ for $y>0$,
\[
 y_h''(t_0)=\frac{(p-1)U}J\Bigl(\frac L{y_f}-\frac M{y_g}\Bigr)
 =\frac{(p-1)U\,y_h(L+M)}{J\,y_fy_g}>0,
\]
using $Ly_g-My_f=y_h(L+M)$ and $L+M=R+B>0$. Thus $y_h$ has a strict local
minimum at $t_0$, so $y_h(t)>y_h(t_0)$ on both punctured sides near $t_0$,
contradicting \eqref{eq:min_hbelow} on the smaller side.

The exchange symmetry transfers (T2) to $\sigma^*=g$. This completes type~T.
\end{proof}

\subsection{Proof of the four-edge lower bound}

\begin{proof}[Proof of Theorem~\ref{thm:fouredge}]
By Lemma~\ref{lem:reductions} we may assume that $\Gamma$ is loopless and
weakly connected with one commodity per endpoint pair. The weak comparison
theorem gives the existence of the one-sided derivatives of $E$ and the weak
inequalities, so it suffices to exclude slope equality. If the underlying
simple graph is a tree, Lemma~\ref{lem:tree} gives the strict inequality for
expansions and shows that contractions do not occur. Otherwise
Lemma~\ref{lem:graphs} applies. If $\mathcal K_{R_s}=\{0\}$,
Lemma~\ref{lem:dimension} excludes slope equality. In the remaining case
$\Gamma$ is a transitive triangle with one duplicated side,
$\dim\mathcal K_{R_s}=1$ and $\mathcal K_{R_l}=\mathcal K(\Gamma)$, and the
smaller regime has type~P or type~T by Lemma~\ref{lem:routespaces}.
Lemmas~\ref{lem:typeP} and~\ref{lem:typeT} exclude slope equality. This proves
the claims for $E$.

For $P=E/S$: $S$ is $C^1$ by Proposition~\ref{prop:Sregular}, and $S(t)>0$
for $t>0$ because every commodity has positive demand and every cost is at
least $\min_ea_e>0$. The quotient rule gives
$P'(t_0^\pm)=E'(t_0^\pm)/S(t_0)-E(t_0)S'(t_0)/S(t_0)^2$, so
\[
 P'(t_0^-)-P'(t_0^+)=\frac{E'(t_0^-)-E'(t_0^+)}{S(t_0)},
\]
which has the same strict sign as the jump of $E$.
\end{proof}

\begin{table}[ht]
\centering\footnotesize
\begin{tabular}{@{}p{3.1cm}p{1.2cm}p{0.5cm}p{1.3cm}p{3.6cm}p{3.6cm}@{}}
\hline
graph or type & $\dim\mathcal K_{R_s},$ $\dim\mathcal K_{R_l}$ & $\sigma^*$ & old long route & expansion & contraction\\
\hline
underlying tree (any $n,e$) & any & -- & -- & bank jump (L.~\ref{lem:tree}) & cannot occur (L.~\ref{lem:tree})\\
$n\ge4$, $e\le4$, not a tree & 0, 1 & -- & -- & dimension lemma & dimension lemma\\
$n=3$, $e\le3$, not a tree & 0, 1 & -- & -- & dimension lemma & dimension lemma\\
$n=3$, $e=4$, not a tree & 0, 1--2 & -- & -- & dimension lemma & dimension lemma\\
$n=3$, $e=4$, directed cycle & 1, 2 & -- & -- & excluded (L.~\ref{lem:graphs}) & excluded (L.~\ref{lem:graphs})\\
\hline
type P & 1, 2 & $f$ & direct & $G''<0$ (P1) & $G''<0$ (P1); degenerate: $y_f<T_0$ (D1)\\
type P & 1, 2 & $f$ & indirect & impossible: bank inactive & impossible: bank inactive\\
type P & 1, 2 & $g$ & indirect & $G'<0$ (P2) & $G'<0$ (P2); degenerate: $G'(t_0^+)<0$ (D2)\\
type P & 1, 2 & $g$ & direct & $G''>0$ and old continuation (P3) & as for expansion (P3); degenerate: tangency and left-side bound (D3)\\
type P & 1, 2 & $h$ & either & exchange symmetry & exchange symmetry\\
type P, equal bank offsets & 1, 2 & $f,g$ & direct & $G\equiv0$ & $G\equiv0$\\
\hline
type T & 1, 2 & $f$ & both & $y_f'>0$, or $A=B=0$ and L.~\ref{lem:monotone} (T1) & same (T1)\\
type T, $y_f,y_g>0$ & 1, 2 & $h$ & both & $y_h''>0$ (T2c) & same (T2c)\\
type T, $y_g(t_0)=0$ & 1, 2 & $h$ & both & $y_h'>0$ (T2a) & same (T2a)\\
type T, $y_f(t_0)=0$ & 1, 2 & $h$ & both & $M=0$, $\Theta$-monotonicity (T2b) & same (T2b)\\
type T & 1, 2 & $g$ & both & exchange symmetry & exchange symmetry\\
\hline
\end{tabular}
\caption{Closing argument for each case of Theorem~\ref{thm:fouredge}.
In type~T the duplicated edge $e_1$ is always loaded at $t_0$; a degenerate
bank in type~P (an edge entering exactly at $t_0$) occurs only for
contractions.}
\label{tab:minimality}
\end{table}

\subsection{Edge-specific exponents}

We end with costs $c_e(z)=a_e+b_ez^{p_e}$, $a_e,b_e>0$, where the exponents
$p_e>1$ may differ between edges. The dimension lemma then fails, because
$Dq$ is no longer a multiple of $c_0-a$. A three-edge transitive triangle
with $c_f=1+z^4$ on $s\to v$, $c_h=\tfrac13+z^2$ on $s\to u$,
$c_g=\tfrac13+\tfrac13z^6$ on $u\to v$ and unit coefficients $R=A=B=1$ has an
equal-slope contraction at $t=1$; along the old continuation, the excess cost
of the indirect route over the direct one is $(t^2-1)^3/3$. For expansions the
triangle is rigid. We use the following form of Descartes' rule of signs.

\begin{lemma}[Descartes' rule for real exponents]\label{lem:descartes}
Let $\varphi(t)=\sum_{i=1}^mc_it^{e_i}$ with real exponents
$e_1<\dots<e_m$ and nonzero real coefficients, and let $V$ be the number of
sign changes in $(c_1,\dots,c_m)$. Then $\varphi$ has at most $V$ zeros in
$(0,\infty)$, counted with multiplicity.
\end{lemma}

\begin{proof}
Induction on $V$. If $V=0$, all terms have one sign and $\varphi$ has no
positive zero. If $V\ge1$, pick $j$ with $c_jc_{j+1}<0$ and
$e^*\in(e_j,e_{j+1})$. The function $\psi=(t^{-e^*}\varphi)'=
\sum_ic_i(e_i-e^*)t^{e_i-e^*-1}$ has coefficients whose signs are those of
$c_i$ reversed for $i\le j$ and unchanged for $i>j$, so it has exactly
$V-1$ sign changes. If $\varphi$ has $Z$ positive zeros with multiplicity,
so has $t^{-e^*}\varphi$, and by Rolle's theorem $\psi$ has at least $Z-1$.
By induction $Z-1\le V-1$.
\end{proof}

\begin{proposition}[Heterogeneous triangle expansions are strict]\label{prop:hetero}
Let $\Gamma$ be the three-edge transitive triangle with costs
$c_\sigma=a_\sigma+b_\sigma z^{p_\sigma}$, $a_\sigma,b_\sigma>0$,
$p_\sigma>1$, and any commodities. At a nested expansion,
$E'(t_0^-)>E'(t_0^+)$.
\end{proposition}

\begin{proof}
Here $\dim\mathcal K(\Gamma)=1$, so $\mathcal K_{R_s}=\{0\}$ by the dimension
lemma's first part, which does not use common exponents. As in its proof, on
the left $z(t)=tq$ with one route per commodity, and the new route is an
$s\to v$ route, so $R>0$. Suppose slope equality holds. Let $G(t)$ be the new
$s\to v$ route's cost minus the old one's along $z=tq$, now for all $t>0$. By
Lemma~\ref{lem:tangencycomparison}, $G>0$ on a left neighbourhood of $t_0$,
$G(t_0)=0$ and $G'(t_0)=0$.

If the old route is indirect,
$G(t)=a_f-c_h((A+R)t)-c_g((B+R)t)$ is strictly decreasing, so $G'(t_0)<0$.
If the old route is direct,
\[
 G(t)=\alpha+\nu_ht^{p_h}+\nu_gt^{p_g}-\beta t^{p_f},\qquad
 \alpha=a_h+a_g-a_f,\ \nu_h=b_hA^{p_h},\ \nu_g=b_gB^{p_g},\
 \beta=b_fR^{p_f}>0 .
\]
After merging equal exponents and dropping zero coefficients this is a
generalised polynomial whose coefficient sequence has at most three sign
changes, and $G\not\equiv0$ because $G>0$ to the left. Let $\mu\ge2$ be the
multiplicity of the zero $t_0$, which is finite because $G$ is analytic on
$(0,\infty)$. If $\mu$ is even, $G$ does not change sign at $t_0$, so $G>0$ on
a punctured right neighbourhood. Then the load $tq$ with the old routes is an
equilibrium to the right, since every unused route has positive excess cost
$G(t)$; by uniqueness it is the equilibrium, and its complete minimum-route
set is $R_s$, contradicting $R_+=R_l$. If $\mu$ is odd, then $\mu=3$ and
Lemma~\ref{lem:descartes} forces exactly three sign changes and no other
positive zero. Three sign changes among at most four terms with a single
negative coefficient $-\beta$ require four distinct exponents and the
pattern $(-,+,-,+)$, so $\alpha<0$ and $G(t)\to\alpha<0$ as $t\to0^+$. With no
zero in $(0,t_0)$, $G<0$ on $(0,t_0)$, contradicting $G>0$ to the left.
\end{proof}

\begin{remark}[Edge counts for edge-specific exponents]\label{rem:hetero}
With edge-specific exponents, Lemmas~\ref{lem:tangencycomparison}, \ref{lem:reversal},
\ref{lem:reductions}, \ref{lem:bank}(a)--(c), \ref{lem:tree} and
\ref{lem:monotone}, and the first part of Lemma~\ref{lem:dimension}, hold
with the same proofs (in Lemma~\ref{lem:bank}(c) the exponent of the entering
edges replaces $p$). With at most three edges, a loopless connected graph is
either a tree of parallel banks, the transitive triangle, or the directed
triangle. In the directed triangle every endpoint pair has a single route, so
$\mathcal K_R=\{0\}$ always and no nested breakpoint exists. Together with
Proposition~\ref{prop:hetero}, three edges never give equal slopes at an
expansion. Since the common-power five-edge example belongs to the larger
class, the minimum edge count for equal-slope expansions with edge-specific
exponents is four or five; which of the two holds is open. For contractions
the minimum is three: two-edge graphs are trees of banks, where contractions
do not occur, and the triangle above gives an example.
\end{remark}
